% Bibliography entries only; article.tex supplies thebibliography.
% Repository links are immutable at the inspected commit.

\bibitem{Repository}
ProveIt Contributors,
\emph{Polylogarithms and their Arithmetic Bridges}.
ProveIt repository, pinned manuscript and editorial-ledger snapshot,
11 October 2026.
\url{https://github.com/VladimirReshetnikov/ProveIt/tree/7bd45777a8a127e5dc69aff8887ca36a79a3059d/Analysis/Polylogarithms/docs/manuscript}.

\bibitem{IncomingParity}
ProveIt research collection,
\emph{Harmonic Parity and Resolvent Identities:
Exact Lerch generators, Stieltjes primitives, and the quartic and higher
Tornheim layers}.
Research report, 11 October 2026; in particular, the centered harmonic
calculus, gap-deletion formulas, and research questions R1, R7, and R8.
Incoming archive at the revision of \cite{Repository}:
\url{https://github.com/VladimirReshetnikov/ProveIt/blob/7bd45777a8a127e5dc69aff8887ca36a79a3059d/docs/incoming/ProveIt_Harmonic_Parity_and_Resolvent_Identities_2026-10-11.zip}.

\bibitem{IncomingIndependent}
ProveIt research collection,
\emph{Independent Orders and Exact Reductions in Polylogarithm--Zeta
Calculus}.
Research report, 11 October 2026; in particular, the entire relative
Hurwitz interpolant and its nonpositive-index reductions.
Incoming archive at the revision of \cite{Repository}:
\url{https://github.com/VladimirReshetnikov/ProveIt/blob/7bd45777a8a127e5dc69aff8887ca36a79a3059d/docs/incoming/ProveIt_Independent_Orders_2026-10-11.zip}.

\bibitem{IncomingHigherReflection}
ProveIt research collection,
\emph{Higher Reflection Identities and Resonant Zeta Calculus}.
Research report, 11 October 2026; in particular, the centered harmonic
parity and resonant-continuation conventions.
Incoming archive at the revision of \cite{Repository}:
\url{https://github.com/VladimirReshetnikov/ProveIt/blob/7bd45777a8a127e5dc69aff8887ca36a79a3059d/docs/incoming/ProveIt_Higher_Reflection_Identities_2026-10-11.zip}.

\bibitem{INY}
Kentaro Ihara, Yayoi Nakamura, and Shuji Yamamoto,
\emph{Interpolant of truncated multiple zeta functions}.
The Ramanujan Journal \textbf{67} (2025), article 1.
\url{https://doi.org/10.1007/s11139-025-01045-2}.
Author preprint: \url{https://arxiv.org/abs/2407.20509}.
Especially Proposition 2.1, Proposition 3.3, Theorem 4.1,
and Corollary 4.2.

\bibitem{Bradley2007}
David M. Bradley,
\emph{A signed analog of Euler's reduction formula for the double zeta
function}.
Author preprint, arXiv:0707.4486 (2007).
\url{https://arxiv.org/abs/0707.4486}.
The unsigned specialization is the classical odd-weight Euler reduction;
see equations (2)--(3) and Proposition 1.

\bibitem{DLMFHurwitz}
National Institute of Standards and Technology,
\emph{NIST Digital Library of Mathematical Functions},
\S25.11, Hurwitz zeta function.
In particular, the shift law, polygamma and Bernoulli special values, parameter
differentiation, and Lerch's log-Gamma identity in
equations 25.11.3, 25.11.12, 25.11.14, 25.11.17, and 25.11.18;
the asymptotic expansion is 25.11.43.
Accessed 11 October 2026.
\url{https://dlmf.nist.gov/25.11}.

\bibitem{DLMFGamma}
National Institute of Standards and Technology,
\emph{NIST Digital Library of Mathematical Functions},
Chapter 5, Gamma function; especially \S5.2 (definitions),
\S5.7 (series expansions), \S5.11 (asymptotic expansions),
and \S5.15 (polygamma functions).
Accessed 11 October 2026.
\url{https://dlmf.nist.gov/5.2};
\url{https://dlmf.nist.gov/5.7};
\url{https://dlmf.nist.gov/5.11};
\url{https://dlmf.nist.gov/5.15}.

\bibitem{DLMFBernoulli}
National Institute of Standards and Technology,
\emph{NIST Digital Library of Mathematical Functions},
\S24.4, Basic properties of Bernoulli and Euler polynomials.
In particular, the difference, reflection, power-sum, and finite-expansion
identities in equations 24.4.1, 24.4.3, 24.4.9, and 24.4.12.
Accessed 11 October 2026.
\url{https://dlmf.nist.gov/24.4}.

\bibitem{DLMFPolylog}
National Institute of Standards and Technology,
\emph{NIST Digital Library of Mathematical Functions},
\S25.12, Polylogarithms.
Accessed 11 October 2026.
\url{https://dlmf.nist.gov/25.12}.

\bibitem{DLMFLerch}
National Institute of Standards and Technology,
\emph{NIST Digital Library of Mathematical Functions},
\S25.14, Lerch's transcendent.
Accessed 11 October 2026.
\url{https://dlmf.nist.gov/25.14}.

\bibitem{DLMFBarnes}
National Institute of Standards and Technology,
\emph{NIST Digital Library of Mathematical Functions},
\S5.17, Barnes' $G$-function (double gamma function).
In particular, the recurrence, normalization, product, and logarithmic
integral formulas in equations 5.17.1, 5.17.3, and 5.17.4.
Accessed 11 October 2026.
\url{https://dlmf.nist.gov/5.17}.

\bibitem{EspinosaMoll}
Olivier Espinosa and Victor H. Moll,
\emph{A generalized polygamma function}.
Integral Transforms and Special Functions \textbf{15} (2004), 101--115.
Author preprint: \url{https://arxiv.org/abs/math/0305079}.
The balanced negapolygamma normalization supplies the fixed constants
used in the primitive ladder.

\bibitem{Vardi}
Ilan Vardi,
\emph{Determinants of Laplacians and multiple gamma functions}.
SIAM Journal on Mathematical Analysis \textbf{19} (1988), 493--507.
\url{https://doi.org/10.1137/0519035}.
The Hurwitz-zeta derivative identity fixes the Barnes-$G$ normalization
used in the endpoint formula.
